\section{Одежда как personality compression: школьная форма, techwear, maid, witch}

Одежда сжимает социальную роль и характер в один читаемый пакет: школьная форма заявляет возраст, институт, принадлежность к группе ещё до реплики; techwear — компетентность, современность, готовность к действию; maid-костюм — service-роль, которую персонаж либо честно занимает, либо иронически подрывает; witch-костюм — знание, дистанция, связь с чем-то старше текущего момента.

Кики из "Kiki's Delivery Service" показывает, что witch-код не обязан быть тёмным или зловещим, чтобы работать: чёрное платье и метла считываются как "ведьма" мгновенно, но функция костюма в сюжете не мистика, а профессия — Кики зарабатывает доставкой на метле, и костюм визуально утверждает связь с традицией одновременно с обычной, бытовой полезностью. Это полезная модель для Рин: witch-код не обязательно означает "загадочная и далёкая", он может означать "владеющая ремеслом, которое для других выглядит волшебством".

Ключевая ошибка здесь та же, что и с силуэтом (6.1): выбирать одежду по эстетике раньше, чем по функции. Witch-код для Рин работает не потому, что "ведьмы это красиво", а потому что он прямо переводит уже установленные черты: дистанцию через знание и странность (1.1), witch-like барьер (3.1), функцию персонажа как соблазна и хаоса в контролируемых дозах (2.7). Maid-код, для сравнения, был бы прямым противоречием её companion-not-servant ядра (0.4).

Для Рин это значит: одежда должна визуально утверждать witch-регистр — или его современный, techwear-адаптированный аналог, "ведьма терминала", а не "ведьма леса", — ближе по духу к рабочему костюму Кики, чем к декоративному мистицизму, и явно избегать service-кодированных силуэтов.
